\section{Введение}\label{sec:intro}

Алгоритмы кластеризации с различными настройками могут формировать существенно разные группировки одних и тех же объектов. Поэтому прикладная задача включает не только построение разбиений, но и выбор результата, соответствующего цели анализа. Такой выбор особенно затруднён, когда известные метки описывают лишь небольшую часть данных: их недостаточно для надёжного сравнения всех кандидатов, однако они содержат информацию о том, какие различия между объектами значимы для задачи.

Для оценки разбиений используются два источника информации. Внутренние индексы характеризуют структуру групп в пространстве признаков. Например, silhouette сопоставляет расстояния внутри кластера и до соседних кластеров~\citep{rousseeuw1987}. Внешние критерии сравнивают разбиение с целевыми метками. Скорректированный индекс Рэнда (Adjusted Rand Index, ARI) учитывает согласие пар объектов с поправкой на случайное совпадение~\citep{hubert1985}. Высокая внутренняя оценка не гарантирует высокого внешнего качества: компактные группы могут объединять разные классы, а один класс может занимать несколько разделённых областей. В работе метки задают прикладной ориентир.

При малой разметке возникает дополнительная неопределённость. Разбиение, максимизирующее ARI на известных метках, может уступать другому разбиению на остальных объектах. При этом важен не только бюджет разметки, но и её состав. Две выборки одинакового размера могут представлять разные классы и содержать разное количество объектов редких категорий. Под покрытием классов далее понимается доля целевых классов, представленных хотя бы одной известной меткой. Эта характеристика используется для последующего анализа эксперимента. Её истинное значение не предполагается доступным правилу выбора.

Ограниченная внешняя информация уже использовалась для выбора моделей полуобучаемой кластеризации и готовых разбиений~\citep{pourrajabi2014,cobs2017}. Исследовались также способы объединения внешних оценок со структурными характеристиками кластеров~\citep{gertrudes2019}. Поэтому исследовательский вопрос заключается в условиях полезности такого объединения. В частности, необходимо установить, помогает ли учёт изменчивости частичного ARI определить, когда внутреннему индексу следует разрешить изменить выбор внешнего критерия.

Рассматриваемое правило сначала выбирает лидера по известным меткам. Для каждого альтернативного разбиения оценивается бутстрепный разброс разности его внешней оценки и оценки лидера. Кандидаты с небольшим отставанием относительно этого разброса образуют допустимое множество. Окончательный выбор определяется внутренним индексом. Такое правило связывает степень вмешательства геометрии с изменчивостью внешнего сравнения. Однако большой разброс может допустить существенно худшие кандидаты, а малый разброс не исключает ошибки при непредставленных классах. Кроме того, лидер выбирается на тех же метках, что используются для оценивания его отрыва, поэтому бутстрепный допуск не интерпретируется как доверительное множество~\citep{efron2014}.

Цель работы --- определить условия, при которых внутренний индекс снижает или увеличивает потерю внешнего качества при выборе из фиксированного пула. Исследование включает три задачи. Во-первых, сопоставить адаптивный допуск с выбором по отдельным критериям и простыми правилами их сочетания. Во-вторых, раздельно оценить влияние объёма разметки, покрытия классов и состава кандидатов. В-третьих, разобрать случаи смены внешнего лидера: как часто происходит вмешательство, какие альтернативы допускаются и к каким последствиям оно приводит на скрытых метках. Предполагаемый вклад состоит в описании границ применимости этого подхода и причин его ошибок. Возможное превосходство конкретного правила проверяется экспериментально.
